\textbf{Актуальность темы исследования}

В современной информационно-коммуникационной среде подавляющее большинство сетевых взаимодействий требует надежной криптографической защиты. Переход индустрии к микросервисной архитектуре и стандартизация протокола TLS 1.3 привели к тому, что алгоритмы на эллиптических кривых Curve25519 (для обмена ключами) и Ed25519 (для цифровых подписей) стали стандартом де-факто благодаря их безопасности и эффективности \cite{RFC7748}. Платформа .NET, являясь одним из лидеров в корпоративной разработке, активно развивается в сторону высокопроизводительных вычислений (High-Performance Computing). Однако интеграция современной криптографии в .NET-приложения сопряжена с трудностями: использование системных библиотек (написанных на C/Assembly) через механизмы P/Invoke нарушает кроссплатформенность и усложняет развертывание, а полностью управляемые библиотеки зачастую не используют современные возможности процессоров, что становится узким местом (bottleneck) при высоких серверных нагрузках. Таким образом, разработка высокопроизводительной, аппаратно-ускоренной криптографической библиотеки на чистом C\# является весьма актуальной задачей.

\textbf{Степень разработанности проблемы}

На сегодняшний день существует ряд реализаций эллиптической криптографии для экосистемы .NET:

\begin{enumerate}
    \item Обертки над нативными библиотеками (например, Libsodium \cite{Libsodium}). Обеспечивают высокую скорость, однако требуют компиляции под каждую целевую операционную систему и архитектуру (x64, ARM), а вызов неуправляемого кода (P/Invoke overhead) нивелирует преимущество скорости на коротких операциях.
    \item Встроенные средства ОС (\texttt{System.Security.Cryptography}). Ограничены платформой Windows (CNG) или Linux (OpenSSL) и не всегда поддерживают современные стандарты кривых в одинаковом объеме.
    \item Управляемые (managed) реализации, такие как индустриальный стандарт Bouncy Castle \cite{BouncyCastle}. Гарантируют 100\% переносимость, но исторически спроектированы под 32-битные архитектуры (использование представления Radix-25.5), активно аллоцируют объекты в куче и не применяют векторные инструкции (SIMD), что критически ограничивает их производительность.
\end{enumerate}

Таким образом, существует открытая ниша для решения, объединяющего безопасность и переносимость чистого управляемого кода с производительностью уровня системных языков программирования за счет глубокой алгоритмической оптимизации и использования SIMD-инструкций.

\textbf{Проблема исследования}

Проблема заключается в том, что стандартные подходы к реализации арифметики конечных полей в управляемых языках не позволяют достичь максимальной пропускной способности процессора. Возникает противоречие между требованием писать переносимый безопасный код (managed C\#) и необходимостью использовать низкоуровневые аппаратные архитектурные особенности (AVX2) для достижения конкурентоспособной скорости в высоконагруженных криптографических системах.

\textbf{Цель работы}

Целью выпускной квалификационной работы является проектирование и реализация высокопроизводительной кроссплатформенной библиотеки криптографических примитивов X25519 и Ed25519 для платформы .NET, обеспечивающей выполнение вычислений за константное время с использованием технологий аппаратного ускорения SIMD и паттернов Zero Allocation.

\textbf{Задачи исследования}

Для достижения поставленной цели были сформулированы следующие задачи:
\begin{enumerate}
    \item Изучить математический аппарат кривых Вейерштрасса, Монтгомери и искривленных кривых Эдвардса, а также алгоритмы арифметики конечного поля $GF(2^{255}-19)$.
    \item Спроектировать математическое ядро с использованием избыточного представления чисел Radix-51 для минимизации операций нормализации и устранения ветвлений.
    \item Разработать архитектуру библиотеки без выделения памяти в управляемой куче с применением современных возможностей языка C\# (\texttt{Span<T>}, \texttt{ref struct}).
    \item Реализовать алгоритм пакетной (batched) обработки независимых скалярных умножений с использованием векторных инструкций AVX2 через механизм \texttt{System.Runtime.Intrinsics}.
    \item Провести профилирование, тестирование по стандартам RFC и бенчмаркинг разработанного решения в сравнении с существующими аналогами (Bouncy Castle).
\end{enumerate}

\textbf{Объект и предмет исследования}

\textbf{Объектом исследования} выступают алгоритмы эллиптической криптографии Curve25519 и Ed25519. \\
\textbf{Предметом исследования} являются программные и аппаратные методы оптимизации арифметики конечных полей в среде исполнения .NET.

\textbf{Методы исследования}

В работе применялись методы математического моделирования (модульная арифметика, проективная геометрия эллиптических кривых), методы структурного и объектно-ориентированного проектирования программного обеспечения, а также эмпирические методы: профилирование производительности, юнит-тестирование по тестовым векторам (Test Vectors) и микро-бенчмаркинг.

\textbf{Научная новизна и практическая значимость}

Научная новизна работы заключается в адаптации подходов аппаратных реализаций (FPGA) к управляемой среде выполнения. Впервые продемонстрирован подход пакетной векторизации (AVX2 Batching) для кривой Curve25519 полностью на языке C\# без использования небезопасного стороннего кода (unsafe C/C++ bindings).

Практическая значимость: разработанная библиотека выложена в открытый доступ (\texttt{github.com/yashelter/dotnet-curve25519-ed25519}), соответствует стандартам RFC 7748 и RFC 8032, и готова к интеграции в реальные .NET-проекты, позволяя увеличить пропускную способность серверов при обработке цифровых подписей и установке защищенных соединений.